% !TeX root = ../../Book2.tex
% 纲要标题：### 6.4 正合性、局部化与标量扩张
% 纲要位置：计划纲要.md，第 1262 行。

\section{正合性、局部化与标量扩张}
\label{b2:sec:0604}

给模添上一组关系，等于取一个商模。张量积会把这些关系送到新的模中，因而商模仍可通过张量计算；但它可能使原来非零的元素变成零，从而破坏单射。两种现象分别解释张量积为何右正合、为何未必正合。局部化与基域扩张也是改变标量的过程，本节还将用张量积描述它们。

% 纲要要点（供目录核对）：
%
% * 张量积的右正合性
% * 有限展示张量后的矩阵计算
% * 张量积不保持单射的实例
% * 局部化、商模与基域扩张中的计算
% * 局部化保持短正合列
%
% ---
%

\subsection{张量积的右正合性}
\label{b2:subsec:0604:01}

\begin{lemma}\label{b2:lem:tensor-cokernel}
设 \(R\) 是含幺环，\(M\) 为幺性右 \(R\) 模，\(N\) 为幺性左 \(R\) 模，\(K\) 为 \(N\) 的子模，\(\iota:K\to N\) 为包含映射。令
\[
H=\operatorname{im}(\operatorname{id}_M\otimes\iota)
\subseteq M\otimes_RN,
\]
即所有 \(m\otimes k\)（\(m\in M,k\in K\)）生成的子群，则自然映射诱导同构
\[
(M\otimes_RN)/H\longrightarrow M\otimes_R(N/K),
\quad m\otimes n+H\longmapsto m\otimes(n+K).
\]
\end{lemma}
\begin{proof}
商映射 \(q:N\to N/K\) 诱导 \(\operatorname{id}_M\otimes q\)，它将 \(H\) 送到零，故给出所述映射。为构造逆，考虑
\[
(m,n+K)\longmapsto m\otimes n+H.
\]
若以 \(n+k\) 代替 \(n\)，像之差为 \(m\otimes k+H=0\)，所以不依赖代表元。双可加性继承自原张量积，而 \(mr\otimes n=m\otimes rn\) 给出平衡性。故此式诱导反向同态。两复合在各自的纯张量或纯张量剩余类上为恒等，所以互逆。
\end{proof}

\begin{definition}\label{b2:def:right-exact-functor}
设 \(\mathcal C,\mathcal D\) 是两个模范畴，\(F:\mathcal C\to\mathcal D\) 是协变加性函子，即 \(F(f+g)=F(f)+F(g)\)。如果 \(\mathcal C\) 中每个正合列
\[
A\longrightarrow B\longrightarrow C\longrightarrow0
\]
在 \(F\) 下仍给出正合列
\[
F(A)\longrightarrow F(B)\longrightarrow F(C)\longrightarrow0,
\]
就称 \(F\) 为\term[youzhenghe-hanzi]{右正合函子}[right exact functor]。这里的模范畴可以分别是不同环的左模或右模范畴。
\end{definition}

\begin{theorem}[张量积的右正合性]\label{b2:thm:tensor-right-exact}
设 \(R\) 是含幺环，\(M\) 为幺性右 \(R\) 模，幺性左 \(R\) 模序列
\[
N'\xrightarrow{f}N\xrightarrow{g}N''\longrightarrow0
\]
正合，则交换群序列
\[
M\otimes_RN'\xrightarrow{\operatorname{id}\otimes f}
M\otimes_RN\xrightarrow{\operatorname{id}\otimes g}
M\otimes_RN''\longrightarrow0
\]
正合。对任意幺性左 \(R\) 模 \(L\)，若幺性右 \(R\) 模序列 \(Q'\xrightarrow{u}Q\xrightarrow{v}Q''\to0\) 正合，则
\[
Q'\otimes_RL\xrightarrow{u\otimes\operatorname{id}_L}
Q\otimes_RL\xrightarrow{v\otimes\operatorname{id}_L}
Q''\otimes_RL\longrightarrow0
\]
也正合。因此 \(M\otimes_R-:R\text{-}\mathbf{Mod}\to\mathbf{Ab}\) 和 \(-\otimes_RL:\mathbf{Mod}\text{-}R\to\mathbf{Ab}\) 都是右正合函子。

若 \(S,T\) 是含幺环，\(M\) 为 \((S,R)\) 双模，\(N',N,N''\) 为 \((R,T)\) 双模且 \(f,g\) 为双模同态，则第一个序列在 \((S,T)\) 双模范畴中正合。若 \(Q',Q,Q''\) 为 \((S,R)\) 双模且 \(u,v\) 为双模同态，\(L\) 为 \((R,T)\) 双模，则第二个序列也在 \((S,T)\) 双模范畴中正合。
\end{theorem}
\begin{proof}
令 \(K=\operatorname{im}f=\ker g\)。由于 \(g\) 满射，\cref{b2:thm:module-first-isomorphism}给出 \(N/K\cong N''\)。\cref{b2:lem:tensor-cokernel}说明 \(\operatorname{id}\otimes g\) 的核恰为所有 \(m\otimes k\)、\(k\in K\) 生成的子群。对每个这样的 \(k\)，存在 \(n'\in N'\) 使 \(f(n')=k\)，故
\[
m\otimes k=(\operatorname{id}\otimes f)(m\otimes n').
\]
反过来，\(\operatorname{id}\otimes f\) 的每个纯张量像均为这种形式，于是其像恰为上述核。又因 \(g\) 满射，每个 \(m\otimes n''\) 都有原像 \(m\otimes n\)，纯张量生成目标，所以 \(\operatorname{id}\otimes g\) 满射。

第一变量的结论可将已证结论用于 \(R^{\mathrm{op}}\)，再用已证明的交换因子同构 \eqref{b2:eq:tensor-opposite-swap}转换回来：它将 \(Q'\otimes_RL\to Q\otimes_RL\) 对应于 \(L\otimes_{R^{\mathrm{op}}}Q'\to L\otimes_{R^{\mathrm{op}}}Q\)，因此第二变量的结论给出原来第一变量的正合性。\cref{b2:prop:tensor-maps}已说明这两个张量函子保持同态的加法。外侧作用分别由 \(s(x\otimes y)=(sx)\otimes y\) 与 \((x\otimes y)t=x\otimes(yt)\) 给出；相应线性条件保证张量后的箭头保留作用。模的核与像在底交换群上计算，故上述序列在所述双模范畴中正合。
\end{proof}

“右正合”中的“右”指序列末端以零结束，不指右模。在定理的记号下，\(\operatorname{coker}f=N/\operatorname{im}f\)，因而上述结论也写成
\[
M\otimes_R\operatorname{coker}f
\cong\operatorname{coker}(\operatorname{id}_M\otimes f).
\]
张量积保留的是由关系生成的商；它没有断言 \(\operatorname{id}\otimes f\) 单射，即使原来的 \(f\) 已经单射。只有左侧或右侧外侧作用时，可分别在定理中取 \(T=\mathbb Z\) 或 \(S=\mathbb Z\)，利用每个交换群的自然整数作用，得到左模或右模范畴中的相应结论。

\ProseParagraphBreak
对环 \(R\) 的双边理想 \(I\)，若要让左 \(R\) 模 \(N\) 的作用经过商环 \(R/I\)，就必须让 \(I\) 的每个元素作用为零。在张量积中，平衡关系给出 \(\bar1\otimes in=\bar i\otimes n=0\)，因此 \(IN\) 中的有限和都被消去。下面的同构说明，无须再加其他关系；右模的情形相同。

\begin{proposition}[商环上的张量计算]\label{b2:prop:tensor-quotient}
设 \(R\) 是含幺环，\(I\) 为 \(R\) 的双边理想，\(N\) 为幺性左 \(R\) 模，\(M\) 为幺性右 \(R\) 模。存在自然同构
\[
(R/I)\otimes_RN\cong N/IN,
\qquad M\otimes_R(R/I)\cong M/MI,
\]
具体公式分别为 \(\bar r\otimes n\mapsto rn+IN\) 与 \(m\otimes\bar r\mapsto mr+MI\)。所得对象分别带左、右 \(R/I\) 模结构。
\end{proposition}
\begin{proof}
第一式的公式不依赖 \(r\) 的代表元，因为 \(i n\in IN\)；它双可加且平衡，所以诱导同态。反向公式为 \(n+IN\mapsto\bar1\otimes n\)。它良定义：\(IN\) 中的元素是有限和 \(\sum_j i_j n_j\)，而
\[
\bar1\otimes\sum_j i_j n_j
=\sum_j\bar i_j\otimes n_j=0.
\]
两个公式互逆，因为 \(\bar1\otimes rn=\bar r\otimes n\)，且 \(1n=n\)。左 \(R/I\) 线性由公式直接给出。第二式用 \(m+MI\mapsto m\otimes\bar1\) 作逆，\(m_j i_j\otimes\bar1=m_j\otimes\bar i_j=0\) 核对关系；其余复合与右线性同样由给出的公式逐项得到。
\end{proof}

对交换环 \(R\) 的两个理想 \(I,J\)，于是有
\[
(R/I)\otimes_R(R/J)\cong R/(I+J).
\]
确实，上一命题先给出 \((R/J)/I(R/J)\)，而 \(I(R/J)=(I+J)/J\)，再用商模的迭代取商得到结果。特别地，正整数 \(m,n\) 满足
\[
(\mathbb Z/m\mathbb Z)\otimes_{\mathbb Z}(\mathbb Z/n\mathbb Z)
\cong\mathbb Z/\gcd(m,n)\mathbb Z,
\]
因为整数理想之和 \(m\mathbb Z+n\mathbb Z=\gcd(m,n)\mathbb Z\)。
取 \(m=2,n=3\)，便得到本章开头的零张量积，这是互素整数理想之和为整个 \(\mathbb Z\) 的结果。

\subsection{张量积不保持单射的例子}
\label{b2:subsec:0604:02}

取整数 \(n>1\)，序列
\[
0\longrightarrow\mathbb Z\xrightarrow{\times n}\mathbb Z
\longrightarrow\mathbb Z/n\mathbb Z\longrightarrow0
\]
正合。与 \(\mathbb Z/n\mathbb Z\) 张量，利用\cref{b2:thm:tensor-associativity-unit}的单位同构及\cref{b2:prop:tensor-quotient}，得到
\[
\mathbb Z/n\mathbb Z\xrightarrow{0}\mathbb Z/n\mathbb Z
\xrightarrow{\operatorname{id}}\mathbb Z/n\mathbb Z\longrightarrow0.
\]
第一个映射为零，因为 \(n a\otimes\bar b=a\otimes n\bar b=0\)；原先的“乘 \(n\)”在张量后仍是乘同一个整数，但此时 \(n\) 已经作用为零。它的来源非零，所以不再单射。第二个映射在上述识别下为恒等，故序列末端仍正合，完全符合右正合定理。

\ProseParagraphBreak
即使张量后单射仍然成立，也不能只凭原映射是包含就不加检查。一个足够条件是原包含有缩回。若 \(i:N'\to N\)、\(r:N\to N'\) 满足 \(ri=\operatorname{id}\)，则
\[
(\operatorname{id}_M\otimes r)\circ(\operatorname{id}_M\otimes i)
=\operatorname{id}_{M\otimes_RN'}
\]
由\cref{b2:prop:tensor-maps}成立，因此张量后的 \(i\) 仍单射，且仍有缩回。结合右正合性，张量积保持每个分裂短正合列。

\ProseParagraphBreak
非分裂时，单射是否保留取决于用来张量的模。下一章的平坦模就是把“对所有单射都保留单射”作为条件。本节的 \(\mathbb Z/n\mathbb Z\) 展示了该条件并非自动满足，不能用“张量积类似向量空间的乘积”代替核的实际计算。

\subsection{局部化与基域扩张}
\label{b2:subsec:0604:03}

商模公式通过让理想中的元素作用为零来改变标量。局部化则让指定的标量成为可逆元；以下只讨论交换环，此时无须区分标量从左还是从右作用。

\ProseParagraphBreak
把分母变成单位，也会使被分母零化的元素变成零。例如，对 \(R=\mathbb Z\)、\(\Sigma=\{1,2,2^2,\ldots\}\)、\(N=\mathbb Z/2\mathbb Z\)，局部化后 \(2\) 可逆，而 \(2\bar1=0\)，所以 \(\bar1\) 的像必须为零。若只按交叉相乘相等识别分式，\((\bar1,1)\) 与 \((0,1)\) 却不能等价，因为原模中 \(\bar1\ne0\)。因此必须允许交叉乘积之差再被某个分母零化。

\begin{definition}\label{b2:def:localized-module}
设 \(R\) 是交换含幺环，\(\Sigma\subseteq R\) 是含 \(1\) 的乘法封闭集，\(N\) 是幺性 \(R\) 模。在 \(N\times\Sigma\) 上规定关系
\[
(n,s)\sim(n',s')\quad\Longleftrightarrow\quad
\text{存在}\;u\in\Sigma\;\text{使}\;u(s'n-sn')=0.
\]
将 \((n,s)\) 的等价类记为 \(n/s\)。在全体等价类上定义以下加法与环局部化 \(\Sigma^{-1}R\) 的标量作用：
\[
\frac ns+\frac{n'}t=\frac{tn+sn'}{st},\qquad
\frac av\frac ns=\frac{an}{vs}.
\]
所得 \(\Sigma^{-1}R\) 模称为 \(N\) 在 \(\Sigma\) 处的\term[mojubu-hua]{局部化}[localization of a module] \(\Sigma^{-1}N\)。其中 \(\Sigma^{-1}R\) 是第一卷定理13.2.2所构造的交换环局部化。
\end{definition}
\symidx{\(\Sigma^{-1}N\)：模 \(N\) 的局部化}

先验证关系确实是等价关系。反身性与对称性直接成立。若 \(u(tn-sn')=0\)、\(v(wn'-tn'')=0\)，则把第一式乘 \(vw\)，第二式乘 \(us\) 后相加，得到
\(uvt(wn-sn'')=0\)。因 \(uvt\in\Sigma\)，传递性成立。

\ProseParagraphBreak
关系中的额外因子也使运算能容许被分母零化的误差。若 \(u(s'n-sn_1)=0\)、\(v(t'n'-tn_1')=0\)，则两种加法结果的交叉相乘之差为
\[
t t'(s'n-sn_1)+s s'(t'n'-tn_1'),
\]
它被 \(uv\) 零化。对标量乘法，若 \(w(v'a-va')=0\)，模分式的关系如前，则两结果的交叉相乘之差可以写成
\[
v'a(s'n-sn_1)+s(v'a-va')n_1,
\]
它被 \(uw\) 零化，故两项运算都不依赖代表元。三个模分式通分到分母的乘积后，两种相加次序有同一分子，因而加法结合律成立；零元是 \(0/1\)，负元是 \((-n)/s\)。标量作用的结合律则由
\[
\frac av\left(\frac bw\frac ns\right)
=\frac{abn}{vws}
=\left(\frac av\frac bw\right)\frac ns
\]
给出，单位作用及两条分配律同样在通分后化为 \(R\) 与 \(N\) 的对应公理。这就得到所述局部化模。

\ProseParagraphBreak
一个分式何时为零，只需在等价关系中与 \((0,1)\) 比较：
\begin{equation}\label{b2:eq:localization-zero}
\frac ns=0\quad\Longleftrightarrow\quad
\text{存在}\;u\in\Sigma\;\text{使}\;un=0.
\end{equation}
因此典范映射 \(N\to\Sigma^{-1}N\)、\(n\mapsto n/1\) 的核，恰为被某个分母零化的元素；局部化没有再消去其他元素。

\begin{proposition}[局部化的张量描述]\label{b2:prop:localization-tensor}
设 \(R\) 是交换含幺环，\(\Sigma\subseteq R\) 是含 \(1\) 的乘法封闭集，\(N\) 是幺性 \(R\) 模，则存在自然的 \(\Sigma^{-1}R\) 模同构
\[
(\Sigma^{-1}R)\otimes_RN\longrightarrow\Sigma^{-1}N,
\qquad \frac as\otimes n\longmapsto\frac{an}s.
\]
其逆映射为 \(n/s\mapsto(1/s)\otimes n\)。
\end{proposition}
\begin{proof}
正向公式对分式 \(a/s\) 的代表元良定义：若 \(u(ta-sb)=0\)，则 \(u(tan-sbn)=0\)。分式加法与模加法给出双可加性，\((ar)n=a(rn)\) 给出 \(R\) 平衡性，故诱导同态，并且保留 \(\Sigma^{-1}R\) 作用。

若 \(n/s=n'/t\)，取 \(u\in\Sigma\) 使 \(u(tn-sn')=0\)。在张量积中，
\[
\frac1s\otimes n-\frac1t\otimes n'
=\frac1{ust}\otimes u(tn-sn')=0.
\]
于是逆映射不依赖代表元。通分公式说明它保持加法，标量作用则由
\[
\frac1{vs}\otimes an=\frac a{vs}\otimes n
=\frac av\left(\frac1s\otimes n\right)
\]
核对。两复合分别恢复 \(n/s\) 与 \((a/s)\otimes n\)，所以互逆。对模同态 \(f:N\to N'\)，两种构造中对应的映射分别为 \((a/s)\otimes n\mapsto(a/s)\otimes f(n)\) 与 \(n/s\mapsto f(n)/s\)，故同构自然。
\end{proof}

\begin{proposition}\label{b2:prop:localization-exact}
设 \(R\) 是交换含幺环，\(\Sigma\subseteq R\) 是含 \(1\) 的乘法封闭集，幺性 \(R\) 模的短正合列为
\[
0\longrightarrow N'\xrightarrow{f}N\xrightarrow{g}N''\longrightarrow0.
\]
则 \(\Sigma^{-1}R\) 模序列
\[
0\longrightarrow\Sigma^{-1}N'
\xrightarrow{\Sigma^{-1}f}\Sigma^{-1}N
\xrightarrow{\Sigma^{-1}g}\Sigma^{-1}N''\longrightarrow0
\]
仍正合，其中局部化同态的公式为 \((\Sigma^{-1}f)(n'/s)=f(n')/s\)、\((\Sigma^{-1}g)(n/s)=g(n)/s\)。
\end{proposition}
\begin{proof}
对任意 \(R\) 模同态 \(h\)，若 \(u(tn-sn_1)=0\)，则 \(u(th(n)-sh(n_1))=h(u(tn-sn_1))=0\)，所以 \(n/s\mapsto h(n)/s\) 良定义；分式运算说明它保留加法与 \(\Sigma^{-1}R\) 作用。由\cref{b2:prop:localization-tensor}的自然性，序列中这两个同态对应于 \(\operatorname{id}_{\Sigma^{-1}R}\otimes f\)、\(\operatorname{id}_{\Sigma^{-1}R}\otimes g\)。\cref{b2:thm:tensor-right-exact}因此保证末端正合。

还须证明 \(\Sigma^{-1}f\) 单射。若 \(f(n')/s=0\)，由\eqref{b2:eq:localization-zero}，存在 \(u\in\Sigma\) 使 \(uf(n')=f(un')=0\)。原映射 \(f\) 单射，故 \(un'=0\)，再由同一零元判据得 \(n'/s=0\)。于是整个局部化序列短正合。
\end{proof}

局部化不仅保留由关系得到的商，还保留单射；下一章将把这项性质表述为 \(\Sigma^{-1}R\) 的平坦性。

\ProseParagraphBreak
例如，\(\Sigma=\{1,p,p^2,\ldots\}\subseteq\mathbb Z\) 时，\(\Sigma^{-1}\mathbb Z=\mathbb Z[1/p]\)。若 \(N=\mathbb Z/p^a\mathbb Z\)，每个元素被 \(p^a\) 零化，由\eqref{b2:eq:localization-zero}可知局部化为零。若 \(\gcd(p,n)=1\)，乘以 \(p\) 已经是 \(\mathbb Z/n\mathbb Z\) 的自同构，分式 \(\bar a/p^j\) 对应 \(p^{-j}\bar a\)，所以此时局部化仍为 \(\mathbb Z/n\mathbb Z\)。

\ProseParagraphBreak
取所有非零整数作分母，得到 \(\mathbb Q\otimes_{\mathbb Z}N\)。由零元判据，典范映射 \(N\to\mathbb Q\otimes_{\mathbb Z}N\)、\(n\mapsto1\otimes n\) 的核恰为挠子模；但有理化并不只是商去挠子模，还允许剩余元素带任意非零整数分母，形成一个 \(\mathbb Q\) 向量空间。例如 \(N=\mathbb Z\) 本来无挠，而 \(\mathbb Q\otimes_{\mathbb Z}\mathbb Z\cong\mathbb Q\)，并不等于 \(\mathbb Z\)。

\ProseParagraphBreak
局部化把环 \(R\) 换成 \(\Sigma^{-1}R\)。另一种标量变换是把域 \(k\) 扩张到更大的域 \(K\)。取域扩张 \(k\subseteq K\) 与 \(k\) 向量空间 \(V\)，则 \(K\otimes_kV\) 是 \(K\) 向量空间，称为 \(V\) 的\term[jiyukuozhang]{基域扩张}[extension of the ground field]。若 \((v_i)_{i\in I}\) 是 \(V\) 的基，则
\[
(1\otimes v_i)_{i\in I}
\]
是 \(K\otimes_kV\) 的 \(K\) 基。事实上，先用交换因子同构，再用\cref{b2:prop:tensor-free-coordinate}，得到坐标同构 \(K\otimes_kV\cong K^{(I)}\)，其公式正是把 \(\sum_i a_i\otimes v_i\) 送到 \((a_i)\)。所以 \(v\mapsto1\otimes v\) 单射，且
\(\dim_K(K\otimes_kV)=\dim_kV\)。这里比较的是两个不同基域上的维数。

\ProseParagraphBreak
线性映射 \(T:V\to W\) 扩张为 \(\operatorname{id}_K\otimes T\)，它在扩张后的基下保留原矩阵。例如，实线性算子
\[
T=\begin{pmatrix}0&-1\\1&0\end{pmatrix}
\]
没有实特征值；在 \(\mathbb C\otimes_{\mathbb R}\mathbb R^2\cong\mathbb C^2\) 上，同一个矩阵有特征向量 \((1,-i)\)、\((1,i)\)，特征值分别为 \(i,-i\)。两个向量组成的矩阵行列式为 \(2i\ne0\)，故形成复基，使扩张算子对角化。基域扩张保留线性映射的公式，却可以改变多项式的分解情况，进而改变可用的标准形。

\ProseParagraphBreak
基域扩张保留自由模的基；对于有限展示模，则可以直接把关系矩阵的系数作用于新的系数模。这也给出了不要求新系数模平坦的张量计算。

\begin{corollary}\label{b2:cor:tensor-finite-presentation}
设 \(R\) 是交换含幺环，\(M,N\) 是幺性 \(R\) 模，且 \(M\) 有有限展示
\[
R^m\xrightarrow{A}R^n\xrightarrow{\pi}M\longrightarrow0,
\qquad A=(a_{ij})\in M_{n\times m}(R),
\]
其中 \(m,n\) 为非负整数，向量均按列书写。定义 \(R\) 模同态
\[
A_N:N^m\longrightarrow N^n,\qquad
(z_j)_{j=1}^m\longmapsto
\left(\sum_{j=1}^m a_{ij}z_j\right)_{i=1}^n.
\]
若 \(e_1,\ldots,e_n\) 为 \(R^n\) 的标准基，则序列
\[
N^m\xrightarrow{A_N}N^n\longrightarrow M\otimes_RN\longrightarrow0,
\qquad (z_i)\longmapsto\sum_{i=1}^n\pi(e_i)\otimes z_i,
\]
正合，并给出自然的 \(R\) 模同构
\[
N^n/\operatorname{im}A_N\longrightarrow M\otimes_RN,
\qquad (z_i)+\operatorname{im}A_N\longmapsto
\sum_{i=1}^n\pi(e_i)\otimes z_i.
\]
零次直和为零模，空和按零理解。
\end{corollary}
\begin{proof}
对展示在第一变量中张量 \(N\)，由\cref{b2:thm:tensor-right-exact}得到
\[
R^m\otimes_RN\xrightarrow{A\otimes\operatorname{id}_N}
R^n\otimes_RN\xrightarrow{\pi\otimes\operatorname{id}_N}
M\otimes_RN\longrightarrow0.
\]
\cref{b2:prop:tensor-free-coordinate}把前两项分别识别为 \(N^m,N^n\)，其逆公式为 \((z_i)\mapsto\sum_i e_i\otimes z_i\)；交换性保证这些坐标同构保留 \(R\) 作用。若 \(\widetilde e_j\) 为 \(R^m\) 的第 \(j\) 个标准基向量，则
\[
(A\otimes\operatorname{id}_N)(\widetilde e_j\otimes z)
=\sum_{i=1}^n e_i a_{ij}\otimes z
=\sum_{i=1}^n e_i\otimes a_{ij}z.
\]
故第一个映射在坐标下恰为 \(A_N\)，第二个映射恰为所述有限求和。正合性说明后者满射且核为 \(\operatorname{im}A_N\)，再由\cref{b2:thm:module-first-isomorphism}得到商模同构。对任意 \(R\) 模同态 \(h:N\to N'\)，逐坐标施加 \(h\) 与 \(A_N,A_{N'}\) 相容，并把上述有限和送到 \(\sum_i\pi(e_i)\otimes h(z_i)\)，这证明自然性。
\end{proof}
